\chapter*{Preface}
A sphere, a doughnut and a coffee cup are shapes, and so is the set of all rotations of space, the set of all unitary matrices, or the
surface on which modular forms live. This book is about such shapes: how to do calculus on them, how to measure their curvature, how to
recognise that two of them are different, and how to integrate over them. Its hero is the idea, due to Riemann, that a space can be
curved without sitting inside anything else, and that everything one needs to know about it can be read from the inside.

The first four chapters are differential geometry. Smooth manifolds, built from the inverse function theorem; vector fields and the
flows they generate, with Lie's bracket and Frobenius' theorem; differential forms, which turn the theorems of Green, Gauss and Stokes
into a single line and prove Brouwer's fixed point theorem in every dimension; and curvature, from Gauss's Theorema Egregium to the
Gauss--Bonnet theorem, where a local quantity integrates to a global invariant.

The next two chapters are algebraic topology. The fundamental group and covering spaces, where the circle's winding number becomes a
group and the subgroups of free groups turn out to be free; then homology, which counts holes in every dimension and proves that
open sets of $\R^n$ cannot be open sets of $\R^m$ unless $n=m$.

The last four chapters bring geometry back to the subjects of this series. Matrix Lie groups, where the rotations of space are
quaternions in disguise and a full turn is not the same as no turn at all; Haar measure and Weyl's integration formula, which give
the law of the eigenvalues of a random unitary matrix, the bridge to random matrices and the zeros of the zeta function; the
hyperbolic plane, the modular group and the geodesic flow, in which continued fractions become the records of geodesics crossing a tiling; and integral
geometry, where Buffon's needle grows into Crofton's formula and the measure of all lines in the plane.

The book assumes Tome~6 (topology, measure and integration), the complex analysis of Tome~7, and the groups and quaternions of
Tome~2. Each chapter ends with exercises, each with a hint and a complete solution, and comes with a Python notebook in which the ideas
are put to work on a computer: robot arms, numerical integrators, areas of maps, curvature of triangle meshes, the word problem in free
groups, topological data analysis, rotations in computer graphics, random matrices, lattice reduction and stereology.
